\section{Method}
We propose an RL weighting framework that augments a Vision Transformer (ViT) with a continuous, stochastic Actor-Critic routing module.
\subsection{Chunkwise Formulation of Linear Attention}
Standard linear attention leverages the associative property to reduce complexity to $\mathcal{O}(N)$. The output for the $i$-th query is typically computed using a global normalization term:
\begin{equation}
    O_i = \frac{\phi(Q_i) \sum_{j=1}^N \phi(K_j)^\top V_j}{\phi(Q_i) \sum_{j=1}^N \phi(K_j)^\top}
\end{equation}

To enable sequential importance assignment, we reformulate this into a causal, chunkwise process. The input sequence of $N$ tokens is partitioned into $T$ non-overlapping sequential chunks of size $C$. At each step $t \in \{1, \dots, T\}$, the backbone processes the current chunk into local Queries ($Q_t$), Keys ($K_t$), and Values ($V_t$).

Rather than computing a full sequence matrix, we maintain two recurrent memory representations to capture global context. The numerator state $S_t \in \mathbb{R}^{D \times D}$ accumulates the cross-covariance of keys and values, while the denominator state $Z_t \in \mathbb{R}^{D \times 1}$ tracks the normalization factor:
\begin{equation}
    S_t = S_{t-1} + \phi(K_t)^\top V_t
\end{equation}
\begin{equation}
    Z_t = Z_{t-1} + \sum_{i \in \text{chunk } t} \phi(K_{i, t})^\top
\end{equation}

The attention output for the current chunk $t$ is then computed as:
\begin{equation}
    O_t = \frac{\phi(Q_t) S_t}{\phi(Q_t) Z_t}
\end{equation}

By reformulating linear attention into a recurrent, chunkwise execution, we isolate the processing of each image segment into a discrete-time Markov Decision Process (MDP). This structure allows a routing agent to sequentially determine the importance of tokens based on local information, global history, and remaining computational resources.

\subsection{Markov Decision Process Formulation}

By reformulating linear attention into a recurrent, chunkwise execution, we isolate the processing of each image segment into a discrete-time Markov Decision Process (MDP). This structure allows a routing agent to sequentially determine the importance of tokens based on local information, accumulated global memory, and remaining computational resources.

\subsubsection{State Space ($\mathcal{S}$)}
At each chunk step $t$, the agent observes a causal state $s_t$ defined as:
\begin{equation}
s_t = \left[ K_t \parallel \mathrm{Proj}(S_{t-1}) \parallel \mathrm{Proj}(Z_{t-1}) \parallel B_t \right]
\end{equation}

The components are:
\begin{itemize}
    \item \textbf{Local Features ($K_t \in \mathbb{R}^{C \times D}$):} Visual features of the current chunk.
    
    \item \textbf{Memory Projection ($\mathrm{Proj}(S_{t-1}) \in \mathbb{R}^{D}$):} A learned projection of the accumulated memory matrix $S_{t-1} \in \mathbb{R}^{D \times D}$.
    
    \item \textbf{Normalization Projection ($\mathrm{Proj}(Z_{t-1}) \in \mathbb{R}^{D}$):} A projection of the normalization accumulator $Z_{t-1}$, encoding the effective mass of stored tokens.
    
    \item \textbf{Dynamic Budget ($B_t \in [0,1]$):} Remaining computational budget.
\end{itemize}

\subsubsection{Action Space ($\mathcal{A}$)}
For each chunk of size $C$, the Actor $\pi_\theta$ outputs:
\begin{equation}
a_t = \{ w_{1,t}, w_{2,t}, \dots, w_{C,t} \}, \quad w_{i,t} \in (0,1)
\end{equation}

The weights are sampled from a squashed Gaussian:
\begin{equation}
a_t \sim \mathrm{SquashedGaussian}(\mu_\theta(s_t), \sigma_\theta(s_t))
\end{equation}

These weights are independent and do not sum to one, allowing the agent to suppress uninformative tokens, dropping the whole chunk if needed.

\subsubsection{Transition Dynamics}
The transition to $s_{t+1}$ is governed by deterministic updates of the memory and normalization accumulators.

The memory update is:
\begin{equation}
S_t = S_{t-1} + \left( (w_t \odot \phi(K_t))^\top V_t \right)
\end{equation}

The normalization term evolves as:
\begin{equation}
Z_t = Z_{t-1} + \sum_{i=1}^{C} w_{i,t} \, \phi(K_{i,t})^\top ,
\end{equation}
where $\phi(\cdot)$ is the feature map associated with linear attention, and $\odot$ denotes element-wise scaling of token features by their corresponding weights.

The budget follows:
\begin{equation}
B_{t+1} = B_t - \frac{1}{N} \sum_{i=1}^{C} w_{i,t}
\end{equation}

\subsubsection{Reward Function ($\mathcal{R}$)}
A terminal reward $R$ is provided at $t = T$:
\begin{equation}
R = \mathrm{TaskMetric} - \lambda_{\text{penalty}} \cdot \max(0, 1 - B_T)
\end{equation}

where:
\begin{itemize}
    \item \textbf{TaskMetric:} Task-specific performance (e.g., mAP or accuracy).
    \item \textbf{Budget Penalty:} Penalizes exceeding the allocated token budget.
\end{itemize}

This formulation encourages the agent to preserve informative tokens while respecting computational constraints.

\subsection{Squashed Gaussian Actor-Critic Network}
To provide the flexibility to down-weight or entirely omit uninformative chunks, we model the routing policy $\pi_\theta(a_t \mid s_t)$ as a Squashed Gaussian distribution[cite: 10]. This independent weighting avoids the mathematical instability of discrete gates and the competitive constraints of Softmax[cite: 10, 11]. The Actor network computes parameters for an unbounded Normal distribution[cite: 12, 13]:
\begin{equation}
    \mu_\theta = \text{MLP}_\mu(s_t), \quad \sigma_\theta = \text{Softplus}(\text{MLP}_\sigma(s_t)) + \epsilon
\end{equation}

During training, actions are sampled and squashed via a hyperbolic tangent to the valid continuous weighting space[cite: 13]:
\begin{equation}
    w_t = \frac{\tanh(x) + 1}{2}, \quad x \sim \mathcal{N}(\mu_\theta, \sigma_\theta^2)
\end{equation}
By decoupling weights from a global sum, the Actor can assign near-zero weights to every token in a redundant chunk, effectively "muting" its contribution to the global memory $S_t$. To maintain mathematical integrity for PPO, the log-probability is corrected using the Jacobian trace of the Tanh transformation:
\begin{equation}
    \log \pi(w_t) = \log p(x) - \sum \log(1 - \tanh^2(x) + \epsilon)
\end{equation}

\subsection{Pretraining: Scale-Aligned Information Distillation}
Phase 0 initializes the Actor using empirical Fisher Information ($\mathcal{I}_i$) as a saliency oracle[cite: 14, 15]. These scores are normalized to $\tilde{\mathcal{I}}_i \in [0, 1]$[cite: 17, 18]. The Actor is trained via supervised learning to align its mean prediction $w_{\mu}$ with this informational utility[cite: 19]:
\begin{equation}
    \mathcal{L}_{\text{Actor}} = \mathbb{E} \left[ \left( w_{\mu}(s_t) - \tilde{\mathcal{I}}_t \right)^2 \right]
\end{equation}

\subsection{Finetuning: Task-Specific Reward Optimization}
In Phase 1, the Actor is optimized via PPO to maximize a task-specific reward $R$[cite: 19]. The environment enforces a decreasing budget $B_t$[cite: 20]. To prevent memory stagnation when tokens are aggressively down-weighted, we introduce State Dilution, injecting a fraction ($\gamma$) of the global context $G$ back into the memory[cite: 21, 22]:
\begin{equation}
    S_t = S_{t-1} + (w_t \odot K_t)^\top V_t + \gamma \cdot (1 - \bar{w}_t) \cdot G^\top V_{\text{global}}
\end{equation}

During deployment (Phase 2), a hard threshold $\tau$ enables physical token purging[cite: 23, 24]. Surviving tokens retain their continuous scaling factors $w_{\text{infer}}$, ensuring the "volume" of information matches the Actor's trained importance weighting[cite: 25].